\scanpage{17}
\mspage{(15)}

\unsure{f)|g)} On peut associer à un espace topologique $\unsure{X}$ ou un ens. $\frac12$ simplicial $\mathcal X_*$ divers invariants alg. \deleted{\unsure{bien}} ``linéaires'' \uline{plus fins} à priori que le $C^{\bullet\bullet}_{\mathrm{DRpd}}$ (le $=$ pour $C_{\bullet\bullet}^{\mathrm{DRpd}}$, en dualisant\dots).

[P.ex. on peut observer que
\[
C^{\bullet\bullet}_{\mathrm{DRpd}*}(k)\unsure{\simeq_0|\simeq}\Gamma^\bullet\Phi_{*k}\otimes_k\Lambda\Psi_{*k}
\]
se déduit de
\[
\inserted{D^{\bullet\bullet}_{*k}=}\Gamma^\bullet\Phi_{*k}\otimes\Lambda^\bullet\Phi_{*k}
\]
(qui est une alg. diff \inserted{$k$-augm.} à p.d.s. qui est une \uline{résolution de $k$}) et de $T\in\Gamma(D^{10}_{*k})$ comme Coker de la multiplication par $dT$ (i.e. on divise par l'idéal engendré par $dT$), ou encore en bidegré donné, \deleted{$\unsure{p+q=0}$} $(p,q)$
\[
\unsure{C}^{p,q}_{*k}\simeq\operatorname{Ker}\left(\diagram{\begin{tikzcd}[column sep=large]
D_*^{p,\unsure{q+1}}\arrow[r,"{\text{produit par }dT}"'] & D_*^{p,q+2}
\end{tikzcd}}\right)
\]
d'où
\[
C^{p,q}_{\mathrm{DRpd}}(\mathcal X_*,k)\simeq\operatorname{Ker}(D^{p,q+1}(\mathcal X_*,k)-D^{p,q+2}(\mathcal X_*,k))
\]
et on peut considérer les structures \inserted{\unsure{dessus}} sur $C^{\bullet\bullet}_{\mathrm{DRpd}}(\mathcal X_*,k)$ comme déduites de certaines structures (à expliciter\dots) sur $D^{\bullet\bullet}(\mathcal X_*,k)$.

On peut aussi définir ``\uuline{La} structure multiplicative à pd \inserted{cohomologique} du type d'homotopie $\mathcal X_*$'', en \inserted{un sens convenable}, qui permet de reconstituer \deleted{\unsure{le}} \inserted{aussi bien} $C^{\bullet\bullet}_{\mathrm{DRpd}}(\mathcal X_*,k)$ que $D^{\bullet\bullet}_k(\mathcal X_*,k)$ (ou les complexes de De Rham-Sullivan). J'ignore si les

\mspage{16}
